\chapter*{Abstract}
\addcontentsline{toc}{chapter}{Abstract}

Table-LLM-Specialist~\cite{xing2024tablellmspecialist} (``Table-Specialist'' below) fine-tunes a language model for a
single table task without human labels. Each task is paired with a \emph{dual} task whose output can check the first:
for NL-to-SQL, a SQL program and a Spark-Scala program written for the same question must return the same rows on
random samples of the table. Only agreeing pairs become training data. The fine-tuned model then generates the next
round's data, for up to three rounds. This report documents a fourteen-day reproduction of the paper's Llama-3.1-8B
results on one workstation. It also covers four experiments beyond the paper.

\textbf{Reproduction.} The paper's training corpus, one of its five NL-to-SQL test sets (Text2Analysis), its test-case
lists, its checkpoint choice and most training settings are unpublished. We substituted a filtered GitTables corpus of
10,053 tables and rebuilt the test sets. We also wrote the parts of the validation pipeline that the released code lacks.
\begin{itemize}
\item \textbf{Base checkpoint.} Following the repository's README, we first used Llama-3.1-8B base. Under every setting
tried, its vanilla scores missed the paper's (WikiSQL 0.326 against 0.525).
\item \textbf{Instruct checkpoint.} Queried through the repository's own chat call, Llama-3.1-8B-Instruct reproduces
the paper's vanilla scores within 0.015 on WikiSQL, WikiTQ and BIRD. After three rounds of Table-Specialist it
reproduces the paper's fine-tuned scores within 0.032 per test set. Spider is the exception: every model of ours,
vanilla included, scores 0.055--0.121 above the paper there. On the zero-shot mean of four test sets, Instruct
Table-Specialist scores $0.6113\pm0.0135$ against the paper's 0.5983.
\item \textbf{Ablation.} On the base model, Table-Specialist beats both vanilla and fine-tuning without validation at
round~3, by $+0.0647$ and $+0.0492$ on the mean of eight cells (paired bootstrap, $p<0.0001$, every seed positive).
\item \textbf{Rounds.} Validation pays off only over rounds. Unvalidated training wins in round~1 and degrades every
round after it.
\item \textbf{Other tasks.} We approximately reproduce the paper's vanilla scores for NL-to-Scala, NL-to-R and error
detection. Our Table-Specialist scores on all three are higher than the paper's.
\end{itemize}

\textbf{Extensions.} We tested two ways to recover the training data that the validator discards.
\begin{itemize}
\item \textbf{Preference learning} (DPO on validated against rejected programs) consistently harmed accuracy, even
after tuning on dev sets ($-0.0209$, $p<0.0001$).
\item \textbf{Repair-TS} repairs each rejected pair once, using only that program's own execution feedback, then
re-validates it with the paper's unchanged Algorithm~3. Its effect depends on the generator and the language:
  \begin{itemize}
  \item On base NL-to-SQL it gains $+0.0421$ at round~3 ($p<0.0001$, all three seeds). A same-data control shows that
  about half of the round-1 gain is not explained by data volume.
  \item On Instruct NL-to-SQL it gains nothing ($-0.0014$, $p=0.55$).
  \item The Scala adapters trained inside those same Instruct runs gain at every round ($+0.0141$ to $+0.0337$,
  $p<0.0001$).
  \end{itemize}
\item \textbf{Self-consistency} at test time adds 0.016--0.025 everywhere.
\item \textbf{Repair for error detection} (repair-ED) helps significantly only at round~1 ($+0.027$ F1, $p=0.03$); it beats
a same-data control but not a plain second sample (best-of-2, $p=0.71$). At rounds 2 and 3 it is still slightly
ahead ($+0.006$, $+0.011$) but not significantly ($p=0.70$, $0.58$), with seeds split in sign.
\item \textbf{NL-to-R}, run end to end with our SQL--R validator: TS gains $+0.128$ over vanilla (mean of 8,
$p<0.0001$), rounds help, and \textbf{repair-TS-R} beats TS-R at every round on every seed ($+0.011$, $+0.020$,
$+0.024$, all $p<0.0001$; zero-shot mean of 4 at round~3: 0.5043 against 0.4760). No same-data control was run for R.
\end{itemize}
We separate measured observations from interpretations and hypotheses throughout. We also grade the evidence for each
candidate contribution, and list the experiments that would be needed to turn the repair findings into a general claim.
